% 本文件用于毕业论文第二章，可直接作为章节文件引用。

\chapter{相关理论与关键技术}
\label{chap:theory}

\section{本章概述}

多采样率工业过程异常检测涉及变量关系建模、局部时序模式提取、频域结构分析以及长程依赖建模等多个方面，单一方法通常难以同时兼顾这些需求。为此，本文在模型设计中综合引入图卷积网络、时序卷积网络、傅里叶变换与 Transformer 等关键技术。本章围绕上述四类方法展开介绍，重点说明其基本思想、核心公式、主要特点及其在工业时间序列建模中的适用性，为后续章节的方法设计与实验分析提供理论基础。

\section{图卷积网络}

\subsection{图结构表示}

与规则网格数据不同，工业系统中的多变量关系通常具有显著的非欧式结构特征。例如，电压、电流、功率、温度和状态量之间往往存在复杂的物理耦合与控制耦合，这类关系难以用传统的一维或二维卷积直接表示。图结构能够以更加自然的方式描述变量之间的关联，因此适合用于工业多变量数据建模。

一般将图表示为
\begin{equation}
    \mathcal{G} = (\mathcal{V}, \mathcal{E}, \mathbf{A}),
\end{equation}
其中，$\mathcal{V}$ 表示节点集合，$\mathcal{E}$ 表示边集合，$\mathbf{A} \in \mathbb{R}^{N \times N}$ 表示邻接矩阵，$N$ 为节点数。当节点 $i$ 与节点 $j$ 存在连接关系时，$A_{ij} > 0$；否则 $A_{ij} = 0$。若每个节点都带有特征向量，则可定义节点特征矩阵
\begin{equation}
    \mathbf{X} = [\mathbf{x}_1, \mathbf{x}_2, \ldots, \mathbf{x}_N]^\top \in \mathbb{R}^{N \times F},
\end{equation}
其中 $F$ 为节点特征维数。

在工业过程场景中，可以将每个传感器变量视为一个节点，将变量间的统计相关性、物理耦合关系或数据驱动学习得到的依赖关系视为边。这样，复杂多变量系统就被转化为图上的特征传播与表示学习问题。

\subsection{图卷积的基本原理}

图卷积网络（Graph Convolutional Network, GCN）的核心思想是通过节点与其邻居之间的消息传递，实现图结构上的局部特征聚合。对于节点 $i$ 而言，其更新后的表示不仅依赖于自身特征，也依赖于邻居节点的信息。因此，图卷积可以看作一种``基于拓扑结构的加权邻域平均''。

经典 GCN 的一层传播公式可写为
\begin{equation}
    \mathbf{H}^{(l+1)} =
    \sigma\!\left(
        \hat{\mathbf{D}}^{-1/2}
        \hat{\mathbf{A}}
        \hat{\mathbf{D}}^{-1/2}
        \mathbf{H}^{(l)}
        \mathbf{W}^{(l)}
    \right),
\end{equation}
其中，$\mathbf{H}^{(l)}$ 表示第 $l$ 层输入特征，$\mathbf{W}^{(l)}$ 为可学习权重矩阵，$\sigma(\cdot)$ 为非线性激活函数，$\hat{\mathbf{A}} = \mathbf{A} + \mathbf{I}$ 表示加入自环后的邻接矩阵，$\hat{\mathbf{D}}$ 为对应的度矩阵，即
\begin{equation}
    \hat{D}_{ii} = \sum_j \hat{A}_{ij}.
\end{equation}

上述公式包含三个关键步骤。首先，通过加入单位矩阵 $\mathbf{I}$ 为每个节点保留自环，使节点在聚合邻居信息时不会丢失自身表示。其次，通过 $\hat{\mathbf{D}}^{-1/2}\hat{\mathbf{A}}\hat{\mathbf{D}}^{-1/2}$ 对邻接矩阵进行对称归一化，避免高度节点在传播过程中产生过强影响，并提升训练稳定性。最后，通过线性变换和非线性映射实现特征空间上的投影与更新。

\subsection{图卷积网络的特点}

与传统全连接建模相比，图卷积具有以下优势：

\begin{itemize}
    \item \textbf{能够显式利用变量间结构关系。} 图卷积不是孤立地处理每个变量，而是在拓扑约束下完成跨变量信息传播，更适合建模复杂工业系统中的耦合关系。
    \item \textbf{参数共享能力较强。} 同一组卷积权重可以应用于不同节点邻域，有助于减少参数规模并提升模型泛化能力。
    \item \textbf{适合与数据驱动图学习结合。} 当真实拓扑未知或难以准确获取时，可以通过可学习邻接矩阵实现自适应结构建模，提高模型的灵活性。
\end{itemize}

当然，图卷积也存在一定局限。例如，浅层图卷积的感受范围有限，而过深的图卷积又容易出现过平滑现象，使不同节点表示逐渐趋同。因此，在实际应用中常需要与时序模型、注意力机制或动态图学习方法结合使用。

\subsection{图卷积在本文中的作用}

本文所研究的多采样率异常检测任务中，不同变量之间存在复杂而隐式的关联关系，仅依靠逐变量时序建模难以充分利用这些信息。图卷积能够在变量维度上实现信息传播，为低频变量的插补和异常状态的识别提供结构约束。因此，图卷积是本文完成跨变量依赖建模的重要基础。

\section{时序卷积网络}

\subsection{一维卷积与因果卷积}

时序卷积网络（Temporal Convolutional Network, TCN）是将卷积神经网络应用于序列建模的一类方法，其核心思想是在时间轴上滑动卷积核，从局部时间窗口中提取动态模式。对于一维时间序列 $\{x_t\}$，卷积运算可写为
\begin{equation}
    y_t = \sum_{i=0}^{k-1} w_i x_{t-i} + b,
\end{equation}
其中 $k$ 为卷积核大小，$w_i$ 为卷积权重，$b$ 为偏置项。

若要求模型在时间步 $t$ 的输出只能依赖当前及历史输入，而不能访问未来信息，则称为因果卷积（Causal Convolution）。因果卷积可保证时间顺序不被破坏，适合在线预测与实时诊断任务。对于重构或窗口级分析任务，也可采用非因果卷积，即同时利用过去和未来上下文，以获得更充分的局部时序信息。

\subsection{空洞卷积与感受野}

普通一维卷积的感受野随网络深度线性增长，当需要建模较长时间依赖时，通常需要堆叠较多卷积层。为在不显著增加参数量的前提下扩大感受野，TCN 常引入空洞卷积（Dilated Convolution）。其计算形式为
\begin{equation}
    y_t = \sum_{i=0}^{k-1} w_i x_{t-d \cdot i} + b,
\end{equation}
其中 $d$ 为膨胀系数。当 $d > 1$ 时，卷积核会以间隔采样的方式覆盖更长时间范围，从而有效扩大感受野。

对于单层空洞卷积，其感受野长度可表示为
\begin{equation}
    R = 1 + (k - 1)d.
\end{equation}
若堆叠多层卷积并逐层增大膨胀系数，则模型能够在保持计算效率的同时建模长时间跨度依赖关系。这也是 TCN 在长序列分析中优于普通卷积的重要原因之一。

\subsection{残差结构与多尺度建模}

标准 TCN 往往采用残差连接结构，以缓解深层网络训练中的梯度消失问题。若将某一卷积模块的输入表示为 $\mathbf{x}$、非线性变换表示为 $\mathcal{F}(\mathbf{x})$，则残差输出可写为
\begin{equation}
    \mathbf{y} = \mathbf{x} + \mathcal{F}(\mathbf{x}).
\end{equation}
这种结构可以使网络更容易学习恒等映射和局部增量信息，从而提高训练稳定性。

在实际工业异常检测任务中，不同异常可能发生在不同时间尺度上。例如，冲击型故障通常表现为短时尖峰，缓慢退化则表现为长期趋势漂移。因此，采用多尺度卷积核并行提取局部特征，有助于同时捕获短时和长时动态模式，增强模型对复杂异常的适应能力。

\subsection{时序卷积网络的特点}

与循环神经网络相比，TCN 具有以下优势：

\begin{itemize}
    \item \textbf{并行计算能力强。} 卷积运算可以在时间维上并行处理，训练效率通常高于依赖递归结构的模型。
    \item \textbf{梯度传播路径更短。} 通过卷积堆叠与残差连接，模型在建模长程依赖时通常更稳定。
    \item \textbf{易于进行多尺度扩展。} 可通过不同卷积核大小、不同膨胀系数灵活控制感受野。
\end{itemize}

但 TCN 也存在一定局限。其局部卷积本质决定了它更擅长提取邻域模式，对于跨越很长时间跨度的全局依赖关系，仍需借助更强的全局建模机制加以补充。

\subsection{时序卷积网络在本文中的作用}

在本文方法中，时序卷积网络主要用于提取多变量序列在时间维上的局部动态特征，尤其适合捕获短时波动、局部趋势以及不同时间尺度上的异常征兆。它与图卷积在建模视角上互补：图卷积强调变量间关系，TCN 强调时间维局部模式，两者结合能够形成更加完整的时域表示。

\section{傅里叶变换与频域分析}

\subsection{离散傅里叶变换}

时间序列既可以在时域中表示，也可以在频域中表示。时域描述信号随时间的变化过程，频域则刻画信号由哪些频率成分组成。傅里叶变换正是完成时域与频域转换的核心工具。

对于长度为 $N$ 的离散时间序列 $\{x[n]\}_{n=0}^{N-1}$，其离散傅里叶变换（Discrete Fourier Transform, DFT）定义为
\begin{equation}
    X[k] = \sum_{n=0}^{N-1} x[n] e^{-j 2\pi kn / N}, \qquad k = 0,1,\ldots,N-1,
\end{equation}
其中 $j = \sqrt{-1}$ 为虚数单位。对应的逆变换为
\begin{equation}
    x[n] = \frac{1}{N}\sum_{k=0}^{N-1} X[k] e^{j 2\pi kn / N}.
\end{equation}

频谱 $X[k]$ 是复数形式，可以进一步分解为实部、虚部、幅值和相位。其中，幅值 $|X[k]|$ 反映第 $k$ 个频率分量的能量大小，相位 $\angle X[k]$ 则表示其相位偏移。对于工业信号而言，周期振荡、谐波成分和噪声干扰往往在频域中具有更加直观的表现形式。

\subsection{快速傅里叶变换}

直接计算 DFT 的时间复杂度为 $O(N^2)$，当序列长度较大时计算开销较高。快速傅里叶变换（Fast Fourier Transform, FFT）通过利用复指数项的对称性与周期性，将复杂度降为
\begin{equation}
    O(N \log N),
\end{equation}
从而显著提高了频域分析效率。对于实值序列，还可以采用仅保留非冗余频谱部分的实数快速傅里叶变换（rFFT），进一步降低存储与计算成本。

因此，FFT 已成为信号处理和深度学习时序建模中最常用的频域分析工具之一。它不仅适用于传统谱分析，也能够与神经网络结合，用于学习频率重要性、增强关键频段表示或构造频域损失函数。

\subsection{频域分析在异常检测中的意义}

很多工业异常在时域中的表现并不十分显著，但在频域中往往具有更好的可分性。例如，谐波畸变会导致特定频段能量增强，机械振动异常会表现为主频漂移或边频出现，噪声污染则常表现为高频能量异常增加。因此，引入频域分析有助于从另一种视角揭示异常特征。

对于多采样率工业数据而言，频域分析还具有两点额外意义。其一，频域表征有助于增强对周期结构和重复模式的建模能力；其二，频域约束可以作为时域重构的补充，缓解单纯依赖均方误差时容易出现的过度平滑问题。当然，在处理缺失值较多或零填充较强的序列时，频域分析也可能受到频谱泄露影响，因此通常需要结合掩码机制、插补策略或频域注意力方法共同使用。

\subsection{傅里叶变换在本文中的作用}

本文所研究的电源系统及相关工业场景具有明显的周期性与频谱结构特征，部分异常在频域中的可分性优于时域。因而，傅里叶变换为本文构建频域分支、强化谐波与周期特征表达提供了必要基础，是实现时频协同异常检测的重要理论支撑。

\section{Transformer}

\subsection{自注意力机制}

Transformer 是近年来序列建模领域最具代表性的深度学习结构之一，其核心是自注意力机制（Self-Attention）。与卷积和循环结构不同，自注意力能够直接建立序列任意两个位置之间的联系，从而更高效地捕获长距离依赖关系。

设输入序列表示为
\begin{equation}
    \mathbf{H} \in \mathbb{R}^{T \times d},
\end{equation}
其中 $T$ 为序列长度，$d$ 为特征维数。首先通过线性变换得到查询矩阵、键矩阵和值矩阵：
\begin{equation}
    \mathbf{Q} = \mathbf{H}\mathbf{W}_Q,\qquad
    \mathbf{K} = \mathbf{H}\mathbf{W}_K,\qquad
    \mathbf{V} = \mathbf{H}\mathbf{W}_V.
\end{equation}
然后，自注意力输出为
\begin{equation}
    \operatorname{Attention}(\mathbf{Q}, \mathbf{K}, \mathbf{V}) =
    \operatorname{softmax}\!\left(
        \frac{\mathbf{Q}\mathbf{K}^{\top}}{\sqrt{d_k}}
    \right)\mathbf{V},
\end{equation}
其中 $d_k$ 为键向量维度。上述公式表明，模型会根据查询向量与键向量的相似度，为不同位置分配不同权重，再对值向量进行加权求和。

\subsection{多头注意力与位置编码}

为了增强模型从不同子空间提取依赖关系的能力，Transformer 通常采用多头注意力（Multi-Head Attention）机制。若共有 $h$ 个注意力头，则第 $i$ 个头的输出为
\begin{equation}
    \mathrm{head}_i =
    \operatorname{Attention}(\mathbf{Q}_i, \mathbf{K}_i, \mathbf{V}_i),
    \qquad i=1,2,\ldots,h.
\end{equation}
将所有头拼接后，再经线性变换得到最终结果：
\begin{equation}
    \operatorname{MHA}(\mathbf{H}) =
    \operatorname{Concat}(\mathrm{head}_1,\ldots,\mathrm{head}_h)\mathbf{W}_O.
\end{equation}

由于自注意力本身对输入顺序不敏感，Transformer 还需要引入位置编码（Positional Encoding）来提供序列位置信息。位置编码可以采用固定的正弦余弦形式，也可以采用可学习参数形式，其本质都是让模型区分``不同时间步''的相对或绝对位置。

\subsection{前馈网络与编码器结构}

在标准 Transformer 编码器中，多头注意力模块后通常接有逐位置前馈网络（Feed-Forward Network, FFN），并配合残差连接与层归一化共同构成编码层。设输入为 $\mathbf{Z}_{l-1}$，则第 $l$ 层编码器可表示为
\begin{align}
    \mathbf{Z}_l' &= \operatorname{LayerNorm}\!\left(
        \mathbf{Z}_{l-1} + \operatorname{MHA}(\mathbf{Z}_{l-1})
    \right), \\
    \mathbf{Z}_l &= \operatorname{LayerNorm}\!\left(
        \mathbf{Z}_l' + \operatorname{FFN}(\mathbf{Z}_l')
    \right).
\end{align}

前馈网络一般包含两层线性变换与非线性激活，其作用是对每个时间位置进行更充分的特征映射。残差连接能够改善梯度传播，层归一化则有助于稳定训练过程。通过多层编码器堆叠，Transformer 可以在全局范围内学习复杂的时间依赖与模式交互。

\subsection{Transformer 的特点}

Transformer 在时间序列建模中的主要优势体现在以下几个方面：

\begin{itemize}
    \item \textbf{全局依赖建模能力强。} 任意两个时间位置都可以直接建立联系，适合处理长序列依赖关系。
    \item \textbf{表达能力丰富。} 多头机制能够从不同子空间学习多种相关模式，提高特征表示能力。
    \item \textbf{适合与其他模块协同工作。} Transformer 可作为全局建模器，与卷积、图神经网络或频域模块结合，形成混合式时序建模框架。
\end{itemize}

其局限性主要在于标准自注意力的时间复杂度为 $O(T^2)$，当序列较长时会带来较高的显存和计算开销。因此，在工程实践中常需要通过局部注意力、稀疏注意力或与局部特征提取模块结合的方式提高效率。

\subsection{Transformer 在本文中的作用}

本文将 Transformer 作为全局重构模块使用，目的是在完成局部时序特征提取和频域增强后，进一步建模整个时间窗口范围内的长程依赖关系。与图卷积和 TCN 相比，Transformer 更擅长捕捉跨越较长时间跨度的整体演化规律，因此能够为异常评分提供更高质量的重构结果。

\section{本章小结}

本章系统介绍了图卷积网络、时序卷积网络、傅里叶变换与 Transformer 四类关键技术。图卷积为多变量耦合关系建模提供了结构化表达；TCN 适合提取局部时序模式和多尺度动态特征；傅里叶变换为周期性、谐波性异常的识别提供了频域分析工具；Transformer 则能够从全局范围建模长程依赖关系。上述理论共同构成了本文后续多采样率异常检测模型设计的技术基础，为第三章提出的时频协同异常检测方法提供了支撑。
